\chapter{Finite element analysis (FEM): strength analysis on a tablet}

This is the most ``beyond expectations'' part of MobileCAD: it is not a viewer that can only display
results---\textbf{the CalculiX solver and the Gmsh mesher are both bundled in the APK}, so modelling,
meshing, solving and reading the contour plot all happen on the device.

\section{What actually runs on the device}

\begin{center}
\small
\begin{tabularx}{\linewidth}{@{}lXX@{}}
\toprule
\textbf{Component} & \textbf{Bundled in the APK} & \textbf{Notes} \\
\midrule
CalculiX ccx 2.22 & Yes (\file{libccx.so}) & Solver for structural and thermal analysis; it really computes on the device \\
Gmsh 4.13.1 & Yes (\file{libgmsh.so}) & Meshing (volume meshes, local refinement)\\
OndselSolver & Yes & Assembly solver \\
Netgen & No & Only creates the object; it cannot be used for meshing \\
Elmer / Z88 / Mystran & No & Only object definitions; no usable solver \\
VTK Python & No & Post-processing uses the built-in FRD parser and drawing pipeline instead \\
\bottomrule
\end{tabularx}
\end{center}

\section{Object model}

A complete analysis consists of these objects (all in the feature tree):

\begin{itemize}
  \item \ui{Analysis} container: holds the material, constraints, mesh, solver and results;
  \item \ui{Material}: solid material (density, Young's modulus, Poisson's ratio, and so on), fluid
        material, non-linear material and others;
  \item \ui{Boundary conditions and loads}: fixed, displacement, force, pressure, contact, spring, bearing,
        gear, temperature, heat flux, and so on;
  \item \ui{Mesh}: Gmsh mesh, mesh region, Gmsh refinement;
  \item \ui{Solver}: \cmd{FEM\_SolverCalculiX};
  \item \ui{Results}: the CCX results object and the post-processing pipeline.
\end{itemize}

\section{Getting started: an analysis in ten steps}

\begin{enumerate}
  \item \textbf{Model}: create a solid on the ``Part'' tab (a cantilever beam, for example) and save it;
  \item Switch to \menupath{More ▾ → FEM};
  \item \ui{Analysis} (\cmd{FEM\_Analysis}): creates a new analysis container;
  \item Select the solid → \ui{Material} (\cmd{FEM\_MaterialSolid}) → fill in the density, Young's modulus
        and Poisson's ratio;
  \item \ui{Fixed} (\cmd{FEM\_ConstraintFixed}) → select the fixed face;
  \item \ui{Force} (\cmd{FEM\_ConstraintForce}) or \ui{Pressure} → select the loaded face → enter the
        magnitude and direction;
  \item \ui{Gmsh mesh from shape} (\cmd{FEM\_MeshGmshFromShape}) → start with a coarse size (5--10 mm, for
        example) → Apply;
  \item \ui{Solver} (\cmd{FEM\_SolverCalculiX}) → \ui{Run solver} (\cmd{FEM\_SolverRun});
  \item Select the results object → \ui{Show result} (\cmd{FEM\_ResultShow}) → choose a field (von Mises
        stress, displacement, principal stress, temperature, and so on) → set the deformation scale;
  \item When you need slices or curves, create a \ui{Post-processing pipeline} and add filters (cut, clip,
        data along a line, data at a point, calculator).
\end{enumerate}

\section{Materials and boundary conditions}

\subsection{Material}

After choosing a solid material with \ui{Material}, fill in its properties: density, Young's modulus,
Poisson's ratio, thermal expansion coefficient, thermal conductivity, and so on. On a phone you can also
pick a preset material from the material library and fine-tune it.

\subsection{Boundary conditions and loads}

The 16 most common boundary conditions are created directly by the app (the selected faces/edges are added
to References automatically): fixed, displacement, force, pressure, contact, spring, bearing, gear, pulley,
temperature, heat flux, initial temperature, rigid body, transform, plane rotation and fluid boundary. The
rest (tie, self-weight, centrifugal force, section print, the electromagnetic constraints, the various
equations) go through FreeCAD's own Python commands.

\section{Meshing}

\begin{itemize}
  \item \ui{Gmsh mesh}: generated from a shape; you can set the maximum and minimum element size and the
        surface detection angle. In the task panel, \ui{Apply} means ``generate the mesh'';
  \item \ui{Mesh region}, \ui{Mesh group}, \ui{Gmsh refinement}: refine the mesh locally;
  \item \ui{Remesh with Gmsh}: re-meshes an existing mesh;
  \item \ui{Clear mesh}, \ui{Mesh info}, \ui{FEM mesh to mesh}: management and conversion;
  \item \textbf{Netgen is not available}---commands that need Netgen fail on mobile.
\end{itemize}

\begin{tip}[Coarse first, fine later]
A phone's CPU and cooling are limited. Use a coarse 5--10 mm mesh first to check that the boundary
conditions are right; once they are, refine to 1--2 mm. One fine-mesh solve in the wrong direction can
waste ten to twenty minutes.
\end{tip}

\section{Solving}

\begin{itemize}
  \item Select the \ui{Solver} (CalculiX) → \ui{Run solver};
  \item The app writes an \file{.inp} file, calls ccx on the device and reads the \file{.frd} results;
  \item A log and a progress indicator are shown during the solve; for large models, plug in the charger
        and keep the screen on;
  \item If it fails, check that the units are consistent, the material is complete, the constraints are
        sufficient (no unconstrained rigid-body motion) and the mesh is valid.
\end{itemize}

\section{Viewing the results}

\ui{Show result} offers these fields: von Mises stress, maximum/minimum principal stress, displacement,
temperature, shear stress, and so on. You can also set the \textbf{deformation scale} (0 = contour plot
only, 1 = true deformation, >1 = magnified).

Post-processing pipelines and filters: \ui{Pipeline from result}, \ui{Warp filter}, \ui{Clip}, \ui{Cut},
\ui{Data along line}, \ui{Data at point}, \ui{Calculator}, \ui{Branch filter}, plus the ``create function''
sources such as plane, sphere, cylinder and box.

\screenshot[0.72]{tablet-fem.jpg}{FEM results: the tree shows the analysis container, the material, the constraints and the CCX result mesh; the view shows the stress contour plot.}

\begin{warn}[Two limitations of post-processing]
The \ui{Glyph} filter and the \ui{Visualization} command are not available in the current port; results are
rendered by colouring the surface vertices. If you need more elaborate visualisation, export the results
and process them in desktop FreeCAD.
\end{warn}

\begin{practice}[Exercise: cantilever beam]
\begin{enumerate}
  \item Create a 100 × 10 × 10 mm box;
  \item FEM → Analysis → Material (steel: density 7850, E 210 GPa, ν 0.3);
  \item Fix one end face; apply a 100 N downward force to the other end face;
  \item Gmsh mesh size 3 mm;
  \item Solve → look at von Mises and the displacement;
  \item Change the force to 200 N, solve again and watch the stress double.
\end{enumerate}
\end{practice}

\section{Performance and limitations: a summary}

\begin{itemize}
  \item Only the CalculiX solver is supported (Elmer/Z88/Mystran are not available); Netgen is not
        available;
  \item No Glyph/visualization filters; post-processing is mainly contour plots and basic pipelines;
  \item Phone memory is limited: keep the element count within a few hundred thousand; if necessary use a
        symmetric model or shell elements;
  \item The solver runs single-threaded or multi-threaded on the device; the speed depends on how well the
        device dissipates heat.
\end{itemize}
